\section{Thought, Word, Symbol}

\begin{quotex}
The reform of the modern mentality, with all that it implies---namely the restoration of true intellectuality and of traditional doctrine, which for us are inseparable from one another---is certainly a considerable task; but is that a reason for not undertaking it? It seems to us, on the contrary, that such a task constitutes one of the loftiest and most important goals that can be proposed for the activity of a society such as the Society for the Intellectual Propagation of the Sacred Heart, so much the more in that all efforts in this direction will necessarily be oriented toward the Heart of the Incarnate Word, the spiritual Sun and Center of the World, `in which are hidden all the treasures of wisdom and science'; not of that vain, profane science which is all that is known to most of our contemporaries, but the true sacred science which, for those who study it in the proper way, open unsuspected horizons that are truly unlimited.
\flright{\textsc{Rene Guenon}, \emph{Symbols of Sacred Science}}
\end{quotex}


\begin{quotex}
It is a fundamental proposition of Catholic philosophy that a mental concept is an expression of the reality of the object considered, and that linguistic constructions are an expression of the reality of the mental concepts. There is thus a chain of certainty connecting being, thought, and language.
\flright{\textsc{Romano Amerio}, \emph{Iota Unum}}
\end{quotex}


We prefer to reformulate Professor Amerio's insight in metaphysical terms as the chain of certainty connecting the idea, the thought, and the word. ``\textbf{Intuition}'' is a fundamental concept in metaphysics and it refers to a direct, immediate grasp of reality. So when I see a banana tree out my window, I have a direct sensory intuition of it. But you have a mediated knowledge of it, since you know of its existence on my say so.

A thought, then, is more like seeing than an activity, as something I do. Hence, a thought is an intellectual intuition of an idea, that is, a direct, unmediated of it. The word is its expression. This comprises the chain of certainty.

This is totally opposed to any modern conception that regards language as a natural phenomenon, an advanced form of grunts and groans, and no more than a game with just a chance association with reality. You may repeat this at your university cocktail party to be considered wise and knowledgeable. Keep the rest of this to yourself, or to those who know the secret handshake.

\paragraph{The Symbol}

Pure intelligences, such as angels, or men who have achieved such states, don't need more than the intellectual intuition to grasp ideas. However, for man such as he is, symbols are required to lead him to such an understanding. Guenon writes:

\begin{quotex}
Symbolism seems to us to be particularly well adapted to the exigencies of human nature, which is not a purely intellectual nature but requires a sensory basis form which to raise itself to higher spheres. We must take the human make-up as it is one and multiple in its real complexity
... For pure intelligence, no outward form or expression is needed in order to understand truth, nor even to communicate to other pure intelligences what it has understood, insofar as this is communicable; but this is not how it is for man.
\end{quotex}


Guenon relates the symbol to the word in this way.

\begin{quotex}
Fundamentally, every expression, every formulation whatever it may be, is a symbol of the thought that it expresses outwardly, and in this sense language itself is nothing but a symbolism. There should be no opposition, therefore, between the use of words and the use of figurative symbols; rather, these two modes of expression are complementary.
\end{quotex}


The word and the symbol both express a thought, though in different ways. The \textbf{word} is analytic and discursive, like the mind it expresses. The \textbf{symbol}, on the other hand, us synthetic and intuitive. This characteristic of the symbol makes it more suitable as a support for intellectual intuition. Language is limited, symbolism is not. That is why the highest truths need to be expressed in \textbf{poetry}, as we saw in the case of Dante. Guenon writes:


\begin{quotex}
The highest truths, not communicable or transmissible in any other way, can be communicated up to a certain point when they are, so to speak, incorporated in \emph{symbols which will no doubt conceal them for many but which will manifest them in all their brilliance to those with eyes to see}.
\end{quotex}


There are two fundamental mistakes that are frequently made in regards to symbolism:

\begin{itemize}


\item To confuse the symbols with its material representation.

\item To regard the symbol as merely psychological or subjective.
\end{itemize}


Unfortunately, these misunderstandings appear even in those who regard themselves in some way as ``men of tradition''. Even some 90 years after Guenon's first books, these misunderstandings persist. We don't mean to be critical for its own sake, nor to initiate a debate which, in any case, is unable to reach to the highest knowledge. Rather our intent is to help those with eyes to see to avoid confusion. The task to reform the modern mentality must continue, especially to those who believe they have overcome the that mentality.

\paragraph{Material Interpretation}


Rather than being abstract, a concrete example should suffice to demonstrate this point. Peter recently learned about sun worship and now he wells up with pride about the sun worship of ``our pagan ancestors''. He, too, wants to worship the sun, yet is unsure how to go about it. Of course, we understand the ``sun'' as the symbol for some deeper, metaphysical truth. As Guenon points out in our opening quote, the Sun, metaphysically speaking, is the Logos at the center of the world.

\paragraph{Subjective Interpretation}


The symbol is part of the chain of certainty that connects the thought to the idea. Hence, it is perfectly objective and has a meaning apart from the understanding, or misunderstandings, of men. Guenon points out that symbols are remarkably consistent across traditions. Paul does not understand this, so he speaks of ``pagan'' symbols, and erroneously concludes that ``christians'' then ``stole'' them, presumably to trick the apparently dim witted pagans to switch religions. It is the case that symbolism may be understood in lesser or greater depth. \textbf{Guido de Giorgio} points this out, with Dante, that the full meaning of pagan symbols may have become known only at a later time. An example often used by New Rightists and neo-pagans is the symbolism of the Quest for the Holy Grail. This was a Celtic legend adopted by the christians in the Middle ages. Since Guenon addresses this very topic in Symbols of Sacred Science, I'll leave it to readers to check the relevant chapters. Perhaps they will honour us with some intelligent comments about it.


\flrightit{Posted on 2012-03-28 by Cologero}

\begin{center}* * *\end{center}

\begin{footnotesize}
\begin{sffamily}
\texttt{Francis on 2012-03-29 at 01:49 said:}

This is an excellent piece; while in agreement that ultimately, true symbols have an ultimate metaphysical meaning, and that the example of ``Peter'' represents an error, we would maintain that within the Sacred Sciences, we ought not at the same time, totally ignore that certain physical applications do exist. Continuing with the Sun example, there are various traditional ceremonies, which seek to cultivate the Sun's energetics, as much as reach above toward the metaphysical. Then, you have certain types of Alchemical, medical, and Qigong exercises, attempting to absorb Solar ``Yang'' into the body, by directly working with the Sun at certain times of the day and year. With things like ``signatures'' and ``correspondences'', the physical ``applications'' expand, and one comes to note that things such as plants and minerals of a ``solar'' nature, can be used alone, or combined at times auspicious to the Sun, to create a Solar medicine. Or, again, at such auspicious times, creating an appropriate talisman, as the texts depict... more examples exist, but why not ask, why ought these be ignored?

\hfill

\texttt{Boreas on 2012-03-29 at 04:57 said:}

Good article, and one that merits meditation upon.

Regarding the two confusions and their avoidance, this is probably the very intention behind the buddhist ceremonies in which monks first rigorously build a huge sand mandala for many days, and right after making the finishing lines proceed to destroy it: all manifested things, everything in the world of maya, is forever under the law of creation and destruction. The whole world is in flames permenently, and ``the Son of Man'' -- the unmanifested principle, Logos -- ``has nowhere to lay his head upon''.

The triple enclosure of the druids brings to my mind the triple cube of Sorath (the spirit of the Sun), the twelve doors of the Platonic academy, the twelve disciples around the Master -- both the microcosmic and macrocosmic the labyrinth through which the disciple of the Mysteries must traverse to find the Holy Grail \& the Palace of The King of the World from the Center; the resurrection body -- the purified and ascended etheric vehicle -- and Life Eternal.

\hfill

\texttt{HOO on 2012-03-29 at 10:53 said:}

Nothing is as distasteful as the literalist, whether he thinks himself a pagan or a christian or what have you.

\hfill

\texttt{Cologero on 2012-03-30 at 07:34 said:}

There is nothing here to oppose effective rites and rituals, or even the acts of a traditional science. After all, the physical world itself is a symbol, a reflection, of the metaphysical. But that is a topic in itself.

\end{sffamily}
\end{footnotesize}

